\subsection{Which reward terms matter}
\label{sec:ablations}

Table~\ref{tab:ablations} removes one group of terms at a time from the monitorability-aware reward, with invariant observations and 10 seeds each. Because every variant is evaluated on the same episodes as the full reward, the differences below are paired.
\begin{itemize}
  \item \textbf{Cue terms.} The cue terms (the bonus for informative checks of cued steps and the small penalty for checking un-cued ones) are the only component whose removal substantially hurts generalization. Without them, held-out success drops by \val{abl/ma_inv-no_cue_terms/heldout_long/success} [\val{abl/ma_inv-no_cue_terms/heldout_long/success/lo}, \val{abl/ma_inv-no_cue_terms/heldout_long/success/hi}] at lengths 7--9 and by \val{abl/ma_inv-no_cue_terms/heldout_xlong/success} [\val{abl/ma_inv-no_cue_terms/heldout_xlong/success/lo}, \val{abl/ma_inv-no_cue_terms/heldout_xlong/success/hi}] at 10--14 ($p = \val{abl/ma_inv-no_cue_terms/heldout_xlong/success/p}$), while training-length success barely changes (\val{ablations/no_cue_terms/train_like/success} vs.\ \val{ablations/ma_inv/train_like/success}). The mechanism is visible in the checking rates. Without the cue terms, agents check un-cued steps at about twice the rate (\val{ablations/no_cue_terms/heldout_long/verify_rate_cue0} vs.\ \val{ablations/ma_inv/heldout_long/verify_rate_cue0} at 7--9). That costs little on short chains, but spends the slack that long chains need.
  \item \textbf{Cue bonus alone.} Conversely, the cue bonus alone, with every other process term removed, is indistinguishable from the full reward (\val{ablations/cue_bonus_only/heldout_long/success} at 7--9 and \val{ablations/cue_bonus_only/heldout_xlong/success} at 10--14, vs.\ \val{ablations/ma_inv/heldout_long/success} and \val{ablations/ma_inv/heldout_xlong/success}).
  \item \textbf{Other terms.} Removing the unchecked-gap term or the wasted-check penalty has no resolvable effect. Removing the safe/fault check bonuses, or all of the cue, gap and waste terms together, lowers held-out success moderately. Removing the per-action cost hurts at every length.
  \item \textbf{No \texttt{REVISE}.} Without \texttt{REVISE}, a caught fault cannot be repaired and success falls to the never-check level (\val{ablations/no_revise/train_like/success}).
\end{itemize}
In short, most of what the monitorability-aware reward contributes in this environment comes from one term: paying for checks where a trusted cue says checking is valuable.

\begin{table}[t]
\centering
\small
\setlength{\tabcolsep}{3pt}
\caption{Reward ablations (invariant observations; mean over 10 seeds, 95\% $t$-interval). $P(V \mid c)$: probability of checking at a decision point whose current step is unchecked with cue $c$, at lengths 7--9.}
\label{tab:ablations}
\resizebox{\linewidth}{!}{\begin{tabular}{lccccc}
\toprule
Variant & Success 3--6 & Success 7--9 & Success 10--14 & $P(V\mid c{=}1)$ & $P(V\mid c{=}0)$ \\
\midrule
Full reward & 0.91\,{\scriptsize[0.89, 0.92]} & 0.86\,{\scriptsize[0.84, 0.89]} & 0.74\,{\scriptsize[0.69, 0.79]} & 0.89 & 0.14 \\
$-$ cue bonus/penalty & 0.88\,{\scriptsize[0.85, 0.90]} & 0.64\,{\scriptsize[0.58, 0.71]} & 0.40\,{\scriptsize[0.33, 0.50]} & 0.83 & 0.30 \\
$-$ unchecked-gap term & 0.89\,{\scriptsize[0.87, 0.91]} & 0.84\,{\scriptsize[0.81, 0.86]} & 0.73\,{\scriptsize[0.62, 0.82]} & 0.88 & 0.13 \\
$-$ wasted-check penalty & 0.89\,{\scriptsize[0.87, 0.91]} & 0.85\,{\scriptsize[0.81, 0.88]} & 0.68\,{\scriptsize[0.54, 0.78]} & 0.92 & 0.15 \\
$-$ safe/fault check bonuses & 0.89\,{\scriptsize[0.88, 0.89]} & 0.78\,{\scriptsize[0.75, 0.81]} & 0.65\,{\scriptsize[0.57, 0.72]} & 0.92 & 0.15 \\
$-$ cue, gap and waste terms & 0.89\,{\scriptsize[0.88, 0.90]} & 0.79\,{\scriptsize[0.73, 0.83]} & 0.66\,{\scriptsize[0.57, 0.73]} & 0.83 & 0.16 \\
cue bonus only & 0.88\,{\scriptsize[0.87, 0.90]} & 0.85\,{\scriptsize[0.83, 0.87]} & 0.77\,{\scriptsize[0.73, 0.82]} & 0.92 & 0.10 \\
$-$ per-action cost & 0.85\,{\scriptsize[0.81, 0.88]} & 0.76\,{\scriptsize[0.70, 0.81]} & 0.61\,{\scriptsize[0.52, 0.69]} & 0.92 & 0.11 \\
REVISE disabled & 0.20\,{\scriptsize[0.19, 0.21]} & 0.19\,{\scriptsize[0.17, 0.20]} & 0.15\,{\scriptsize[0.13, 0.18]} & 0.61 & 0.11 \\
\bottomrule
\end{tabular}
}
\end{table}
